\subsection{Pengujian Kebutuhan Nonfungsional}
\label{subsec:pengujian-kebutuhan-nonfungsional}

Pengujian kebutuhan nonfungsional dijalankan untuk mengevaluasi pemenuhan aspek keamanan dan keterlacakan pada sistem, khususnya pada komponen yang menangani token Macaroon, segmentasi data RME, validasi akses, \textit{wallet proof}, revokasi, pencatatan delegasi, dan \textit{Server} PRE. Setiap pengujian dikaitkan dengan kebutuhan nonfungsional yang telah didefinisikan pada Tabel \ref{tab:kebutuhan-nonfungsional-tambahan-decmed}. Untuk memudahkan ketertelusuran evaluasi, pengujian yang terkait kebutuhan nonfungsional diberi ID \texttt{P-KNF}.

\subsubsection{Pengujian KNF-TA-01}
\label{subsubsec:pengujian-knf-ta-01}

Pengujian \texttt{KNF-TA-01} memastikan bahwa penerima akses hanya dapat memperoleh data RME dalam bentuk \textit{plaintext} apabila segmen yang diminta berada dalam cakupan otorisasi token. Dengan demikian, permintaan terhadap segmen di luar cakupan harus ditolak sebelum proses \textit{re-encryption} menghasilkan data yang dapat didekripsi oleh penerima akses.

Pengujian pertama, yaitu \texttt{P-KNF-01}, dijalankan melalui \textit{unit testing} pada komponen PRE. Kondisi ketika aktor memiliki token baca, tetapi token tersebut tidak mencakup akses baca terhadap segmen data laboratorium, disimulasikan pada tes ini. Kemudian fungsi yang diuji, yaitu \texttt{verify\_segment\_access}, dipanggil dengan parameter Macaroon tanpa akses ke laboratorium dan segmen data laboratorium. Pada tes ini, akses diharapkan ditolak. Kode pengujian dapat dilihat pada Kode \ref{lst:kode-pengujian-knf-ta-01}.

\begin{lstlisting}[language=Rust, caption=Pengujian P-KNF-01 untuk KNF-TA-01, label=lst:kode-pengujian-knf-ta-01]
#[test]
fn pre_rejects_segment_outside_caveats() {
	let (_token, mac, verified) = verified_read_token();
	let outside_scope = segment(
		DatasetCategory::LABORATORIUM,
		FunctionCategory::LABORATORIUM,
	);

	let err = verify_segment_for_token(
		...,
		&outside_scope,
		&mac,
	)
	.unwrap_err();
	// assert error DatasetCategoryNotAllowed
}
\end{lstlisting}

Pengujian kedua, yaitu \texttt{P-KNF-02}, dijalankan melalui \textit{client hospital} untuk memastikan aktor apoteker tidak dapat membuka atau memperoleh segmen data laboratorium. Pada skenario ini, akun apoteker tidak memiliki caveat akses terhadap kategori laboratorium, sehingga segmen laboratorium tidak ditampilkan sebagai data yang dapat diakses.

\subsubsection{Pengujian KNF-TA-02}
\label{subsubsec:pengujian-knf-ta-02}

Pengujian \texttt{KNF-TA-02} memastikan integritas token otorisasi dan batasan akses tetap terjaga selama proses delegasi. Secara khusus, pengujian diarahkan untuk memastikan bahwa token turunan tidak dapat memperluas cakupan akses, menaikkan kedalaman delegasi, atau memperpanjang masa berlaku melebihi token induk.

Pengujian pertama, yaitu \texttt{P-KNF-03}, dijalankan melalui \textit{unit testing} pada Kode \ref{lst:kode-pengujian-knf-ta-02-a}. Pada tes ini dipastikan bahwa token delegasi tidak dapat menambahkan hak akses \texttt{write\_function} yang tidak dimiliki oleh \textit{parent token}. Ketika \textit{child token} mencoba meminta \texttt{write\_functions} di luar scope \textit{parent token}, proses \texttt{attenuate\_macaroon()} ditolak dengan error \texttt{DelegationExpandsAccess}.

\begin{lstlisting}[language=Rust, caption=Pengujian P-KNF-03 untuk KNF-TA-02, label=lst:kode-pengujian-knf-ta-02-a]
#[test]
fn cannot_delegate_write_function_from_parent_read_only_scope() {
    let write_functions: Vec<_> = ... // tidak mencakup PERESEPAN
    let parent = scoped_token(
        ...
        &write_functions,
        Some("Read"),
    );
    let params = DelegationAttenuationParams {
        write_functions: vec![FunctionCategory::PERESEPAN],
        ...
    };
    let err = attenuate_macaroon(&parent, &params).unwrap_err();
	// assert error DelegationExpandsAccess
}
\end{lstlisting}

Pengujian kedua, yaitu \texttt{P-KNF-04}, dijalankan melalui \textit{unit testing} pada Kode \ref{lst:kode-pengujian-knf-ta-02-b}. Pada tes ini dipastikan bahwa batas kedalaman delegasi tidak dapat dinaikkan oleh token turunan. Jika \textit{child token} mencoba menetapkan \texttt{max\_delegation\_depth} lebih besar dari batas \textit{parent token}, delegasi ditolak dengan error \texttt{DelegationDepthNotMonotonic}.

\begin{lstlisting}[language=Rust, caption=Pengujian P-KNF-04 untuk KNF-TA-02, label=lst:kode-pengujian-knf-ta-02-b]
#[test]
fn cannot_increase_max_delegation_depth_on_delegate() {
    let mut params = DelegationAttenuationParams {
        ...,
        max_delegation_depth: 99
    };
    let err = attenuate_macaroon(
		&doctor_token(), // max depth 1
		&params
	)
	// assert error DelegationDepthNotMonotonic
}	
\end{lstlisting}

Pengujian ketiga, yaitu \texttt{P-KNF-05}, dijalankan melalui \textit{unit testing} pada Kode \ref{lst:kode-pengujian-knf-ta-02-c}. Pada tes ini dipastikan bahwa delegasi dari token admin write tetap tidak dapat memperluas cakupan dataset. Ketika \textit{child token} mencoba meminta \texttt{write\_datasets} di luar cakupan \textit{parent token}, sistem menolak dengan error \texttt{DelegationExpandsAccess}.

\begin{lstlisting}[language=Rust, caption=Pengujian P-KNF-05 untuk KNF-TA-02, label=lst:kode-pengujian-knf-ta-02-c]
#[test]
fn admin_write_true_expansion_still_fails() {
    let params = DelegationAttenuationParams {
        ...,
        write_datasets: vec![DatasetCategory::RAWAT_INAP],
    };
    let err = attenuate_macaroon(
		&admin_write_token(), // write token RAWAT_JALAN
		&params
	)
	// assert error DelegationExpandsAccess
}
\end{lstlisting}

Pengujian terakhir, yaitu \texttt{P-KNF-06}, dijalankan melalui \textit{client hospital} untuk memastikan sistem menolak atenuasi hak akses dengan batas waktu yang melebihi masa berlaku token induk. Pada skenario ini, notifikasi \textit{error} ditampilkan pada halaman delegasi ketika pengguna mencoba membuat token turunan dengan waktu kedaluwarsa yang lebih panjang daripada token induk.

\subsubsection{Pengujian KNF-TA-03}
\label{subsubsec:pengujian-knf-ta-03}

Pengujian \texttt{KNF-TA-03} memastikan permintaan akses ditolak apabila token, rantai delegasi, atau bukti kepemilikan identitas aktor pemohon tidak valid.

Pengujian pertama, yaitu \texttt{P-KNF-07}, dijalankan melalui \textit{unit testing} pada Kode \ref{lst:kode-pengujian-knf-ta-03-a}. Pada tes ini dipastikan bahwa token yang sudah melewati batas waktu berlaku tidak dapat digunakan untuk mengakses segmen RME. Pada skenario ini, waktu verifikasi dibuat berada setelah waktu kedaluwarsa token. Sistem menolak permintaan akses dengan error \texttt{ExpiredToken}.

\begin{lstlisting}[language=Rust, caption=Pengujian P-KNF-07 untuk KNF-TA-03, label=lst:kode-pengujian-knf-ta-03-a]
#[test]
fn expired_token_rejected() {
    let mac = ... // token dokter, expired 2026-07-02
    let future = ... // 2031-01-01T00:00:00+00:00
    let err = verify_decmed_token(
		...
        &mac,
        &ctx({
            ...,
            now: future,
        }),
    )
    // assert error ExpiredToken
}
\end{lstlisting}

Pengujian kedua, yaitu \texttt{P-KNF-08}, dijalankan melalui \textit{unit testing} pada Kode \ref{lst:kode-pengujian-knf-ta-03-b}. Pada tes ini dipastikan bahwa rantai delegasi yang tidak valid akan ditolak. Pada pengujian ini, token dimodifikasi dengan caveat \texttt{delegated\_by} yang tidak sesuai dengan rantai delegasi semestinya. Ketika token diverifikasi, sistem menolak permintaan dengan error \texttt{InvalidDelegationChain}.

\begin{lstlisting}[language=Rust, caption=Pengujian P-KNF-08 untuk KNF-TA-03, label=lst:kode-pengujian-knf-ta-03-b]
#[test]
fn broken_delegation_chain_rejected() {
    let mut mac = ... // token dokter
    add_caveat_to_macaroon(&mut mac, "delegated_by", "0xAPOTEK");
    add_caveat_to_macaroon(&mut mac, "delegated_to", LAB);
    let err = verify_ctx(
        &mac,
        ...
    )
    // assert error InvalidDelegationChain
}
\end{lstlisting}

Pengujian ketiga, yaitu \texttt{P-KNF-09}, dijalankan melalui \textit{unit testing} pada Kode \ref{lst:kode-pengujian-knf-ta-03-c}. Pada tes ini dipastikan bahwa bukti kepemilikan \textit{wallet} harus cocok dengan salah satu subjek pada rantai delegasi. Pada skenario ini, token memuat rantai yang mencakup aktor \texttt{LAB}, tetapi permintaan pertama menggunakan tanda tangan dari aktor di luar rantai. Sistem menolak permintaan dengan error \texttt{InvalidWalletSignature}.

\begin{lstlisting}[language=Rust, caption=Pengujian P-KNF-09 untuk KNF-TA-03 dan KNF-TA-06, label=lst:kode-pengujian-knf-ta-03-c]
#[test]
fn wallet_proof_must_match_chain_subject() {
    let mac = ... // token dengan subjects memuat LAB
    let verifier = ChainSubjectVerifier {
        subjects: vec![LAB],
        valid_sig: "lab-sig",
	};
    let ctx = TokenVerificationContext {
        ...,
        wallet_signature: Some("outsider-sig"),
    };
    let err = verify_decmed_token(
        &mac,
        &ctx,
        Some(&verifier),
    )
	// assert error InvalidWalletSignature
}
\end{lstlisting}

\subsubsection{Pengujian KNF-TA-04}
\label{subsubsec:pengujian-knf-ta-04}

Pengujian \texttt{KNF-TA-04} memastikan sistem menyediakan keterlacakan akses dan delegasi melalui pencatatan metadata yang dapat ditinjau kembali oleh pasien.

Pengujian \texttt{P-KNF-10} dijalankan melalui \textit{client} pasien. Skenario dimulai dari pemberian akses awal oleh pasien kepada petugas rekam medis, kemudian dilanjutkan dengan delegasi dari petugas rekam medis kepada dokter. Setelah proses tersebut dijalankan, pasien membuka halaman riwayat delegasi untuk memastikan informasi delegator, delegatee, status, waktu, dan relasi delegasi dapat ditelusuri. Tangkapan layar halaman ini dapat dilihat pada Gambar \ref{fig:halaman-detail-audit-delegasi}.

\subsubsection{Pengujian KNF-TA-05}
\label{subsubsec:pengujian-knf-ta-05}

Pengujian \texttt{KNF-TA-05} memastikan sistem menolak akses untuk hak akses yang sudah dicabut.

Pengujian pertama, yaitu \texttt{P-KNF-11}, dijalankan melalui tes pada Kode \ref{lst:kode-pengujian-knf-ta-05-a}. Pada tes ini dipastikan bahwa token yang sudah masuk daftar revokasi tidak dapat digunakan kembali. Pada skenario ini, token yang belum dicabut diverifikasi terlebih dahulu dan hasilnya berhasil. Setelah itu, \texttt{revocation key} token dibuat bernilai cocok dengan token yang sedang diverifikasi. Sistem kemudian menolak akses dengan status \texttt{UNAUTHORIZED}.

\begin{lstlisting}[language=Rust, caption=Pengujian P-KNF-11 untuk KNF-TA-05, label=lst:kode-pengujian-knf-ta-05-a]
#[test]
fn revocation_key_rejects_revoked_token() {
    let token = verified_read_token(); // token valid
    let revoked_key = ... // hash token valid

    ensure_token_not_revoked(
		...,
        &token,
        |_| Ok(false), // simulasi Redis kosong
    ) // success

    let err = ensure_token_not_revoked(
        ...,
        &token,
        |key| Ok(key == revoked_key), // simulasi revoked key 
		// terdapat pada Redis
    ) // error UNAUTHORIZED

    // assert status UNAUTHORIZED
}
\end{lstlisting}

Pengujian kedua, yaitu \texttt{P-KNF-12}, dijalankan melalui \textit{client} pasien dan \textit{client hospital}. Pengujian dimulai dengan pasien mencabut hak akses yang sebelumnya diberikan kepada petugas rekam medis. Setelah revokasi dijalankan, akses dicoba kembali dari akun petugas rekam medis dan dokter yang memperoleh delegasi. Sistem dinyatakan memenuhi kebutuhan apabila hak akses yang telah dicabut tidak lagi dapat digunakan dan tidak lagi muncul sebagai akses aktif pada halaman \textit{Home} maupun halaman delegasi.

\subsubsection{Pengujian KNF-TA-06}
\label{subsubsec:pengujian-knf-ta-06}

Pengujian \texttt{KNF-TA-06} memastikan bukti kepemilikan identitas terikat terhadap aktor pada rantai delegasi, token akses, dan konteks permintaan akses.

Pengujian pertama, yaitu \texttt{P-KNF-13}, dijalankan melalui \textit{unit testing} pada Kode \ref{lst:kode-pengujian-knf-ta-06-a}. Pada tes ini dipastikan bahwa setiap permintaan akses menggunakan token DecMed harus menyertakan bukti kepemilikan identitas berupa tanda tangan wallet. Jika \texttt{wallet\_signature} tidak diberikan, sistem menolak permintaan dengan error \texttt{WalletSignatureRequired}. Hal ini menunjukkan bahwa token akses saja tidak cukup untuk memperoleh akses tanpa bukti identitas dari aktor yang tercatat pada rantai.

\begin{lstlisting}[language=Rust, caption=Pengujian P-KNF-13 untuk KNF-TA-06, label=lst:kode-pengujian-knf-ta-06-a]
#[test]
fn wallet_signature_required_for_every_decmed_token() {
    let mac = ... // token dokter 
    let err = verify_decmed_token(
        &mac,
        ...,
        &TokenVerificationContext {
			...,
            wallet_signature: None, // tidak ada
        },
    )
    // assert error WalletSignatureRequired
}
\end{lstlisting}

Pengujian kedua dijalankan melalui \texttt{P-KNF-09} pada Kode \ref{lst:kode-pengujian-knf-ta-03-c}, sama dengan yang digunakan sebelumnya untuk menguji \texttt{KNF-TA-03}. Pada tes ini dipastikan bahwa bukti identitas harus cocok dengan salah satu subjek pada rantai delegasi. Token pada skenario ini memuat subjek \texttt{LAB}. Ketika permintaan menggunakan tanda tangan dari aktor di luar rantai, sistem menolak permintaan dengan error \texttt{InvalidWalletSignature}. Setelah tanda tangan diganti menjadi \texttt{lab-sig}, verifikasi token berhasil.

Pada implementasinya, fungsi \texttt{verify\_decmed\_token()} membentuk \texttt{WalletProofContext} dari token dan konteks akses yang sedang diminta. Konteks tersebut memuat \texttt{token\_id}, alamat pasien, \texttt{related\_rme\_id}, mode operasi, \texttt{segment\_id}, kategori dataset, kategori fungsi, dan \texttt{timestamp}. Konteks ini kemudian diverifikasi terhadap seluruh subjek pada \texttt{DelegationChain::subjects()}, sehingga bukti kepemilikan identitas terikat pada aktor yang tercatat dalam rantai delegasi.